\section{Erkenntnisgewinn und Triangulation in der Fallstudie}

In diesem Abschnitt werden die Erkenntnisse der quantitativen und qualitativen Methoden dargelegt. Davon ausgehend werden die Erkenntnisse beschrieben, die sich aus der Triangulation der Methoden, Daten und theoretischen Perspektiven ergeben. Schließlich werden der Beitrag der Triangulation zur Validität und Zuverlässigkeit der Ergebisse diskutiert.

\subsection{Systematik zur Beschreibung des Erkenntnisgewinns}

Um den Erkenntnisgewinn durch quantitative und qualitative Methodik sowie 
ihre Triangulation systematisch zu beschreiben, wird eine Verankerung im 
Forschungszweck und im Forschungsinteresse der Studie vorgenommen. Als 
methodologische Grundlage dient die Unterscheidung in  
deskriptives, exploratives und explanatives Erkenntnisinteresse nach Döring und Bortz (2016, Kap. 7.5 TODO). Jede Forschungsfrage der Studie wird dabei einem 
Erkenntnisinteresse zugeordnet und der Methode, die zu ihrer Beantwortung 
eingesetzt wurde. Aus dieser Zuordnung ergibt sich folgende 
studienbezogene Übersicht der Erkenntnistypen (ET1, ET2, ...):

\begin{itemize}
  \item \textbf{ET1 – Wirkung} (deskriptiv): F: Was hat sich durch das Programm „Gute Chancen für alle Kinder“ verändert? Welche Erfolge konnten durch das Programm „Gute Chancen für alle Kinder“ verzeichnet werden?
  \item \textbf{ET2 – Zielerreichung} (deskriptiv), F: Entfalten die umgesetzten Maßnahmen die beabsichtigten Wirkungen?
  \item \textbf{ET3 – Informationszugang} (deskriptiv), F: Wie gut werden Informationen über Hilfen und Angebote für von Armut betroffenen oder bedrohten Familien, Kindern und Jugendlichen zugänglich gemacht?
  \item \textbf{ET4 – Nutzung \& Bekanntheit} (deskriptiv), F: Welche der umgesetzten Maßnahmen kommen bei den Adressat*innen an? Wie bekannt sind die verschiedenen Angebote und Hilfen des Tübinger Präventionskonzepts gegen Kinderarmut bei den Adressat*innen?
  \item \textbf{ET5 – Hürden} (explorativ), F: Welche Hürden bestehen (weiterhin)?
  \item \textbf{ET6 – Bedarfe} (explorativ), F: Welche weiteren spezifischen Bedarfe haben die von Armut betroffenen oder bedrohten Familien, Kinder und Jugendlichen in Tübingen?
\end{itemize}

In ET1 und ET4 wurden jeweils zwei Forschungsfragen zusammengefasst aufgrund einer Nähe im Forschungsinteresse. Es fällt auf, dass keine explanative Forschungsfrage formuliert wurde. Die Studienautoren hatten also keine Hypothese zu einer Kausalbeziehung, die sie testen wollten. Stattdessen liegt das Interesse einerseits auf einer Beschreibung der Situation und andererseits auf der Exploration von Hürden und Bedarfen zur Konzeptweiterentwicklung.

Neben den sechs Erkenntnistypen wird außerdem auf Erkenntnisse eingegangen, die zur Entwicklung nachfolgender Forschungsmethoden beitragen. Da es sich beim Studiendesign weder um eine einzelne Untersuchung noch um parallele oder voneinander unabhängige Teilstudien, sondern um ein sequentielles Mixed-Methods-Design handelt, kommen diesen Erkenntnissen eine besonderer Relevanz zu.

\subsection{Erkenntnisse der quantitativen Befragung}

Die Online-Umfrage unter 355 von Armut betroffenen oder bedrohten Familien aus Tübingen bildet eine statistische Grundlage für die Beantwortung der Foschungsfragen. Die Ergebnisse der Umfrage wurden univariant und bivariant analysiert. Dabei konnten etwa Anspruchnahmen in Relation zu der soziodemographischen Situation der Befragten gestellt werden. So bewerten zum Beispiel Ein-Eltern-Haushalte ihre finanzielle Situation schlechter.

In der Umfrage wird ein gesonderter Fokus auf die Teilhabesituation der Kinder gelegt. Dies geschieht konkret anhand eines Fragebogens zur Deprivation (TODO). Die Erforschung der derzeitigen Situation bietet allerdings nur geringe Erkenntnisgewinne bezüglich der erzielten Wirkung (ET1) und Zielerreichung (ET2). Um eine Wirkung messen zu können wäre ein Vergleich mit der finanziellen Situation und Teilhabesituation vor einigen Jahren nützlich. Wirkung wurde teilweise in den Antworten auf offene Fragen benannt wie etwa, dass die Leistungen die Integration des Kindes erleichterten („Meine jüngste Tochter wurde in der Grundschule […] gut aufgenommen. Heute spricht sie fließend Deutsch. Vielen Dank […] für "Schwimmen für alle Kinder". […] Das hat bei der Integration meiner Tochter sehr gut funktioniert.“ TODO). Inwiefern die Maßnahmen die beabsichtigte Wirkung erzielen (ET2), wird hier nicht erhoben.

Die Umfrage konnte Erkenntnisse dazu liefern, wie gut der Informationszugang für die Zielgruppe ist (ET3). Dafür wurde erhoben, wie informiert sich die Befragten fühlen und auch über welche Quellen und Anlaufstellen sie sich informieren. Ein Mangel an Informationsfluss und besonders der Wunsch nach einer einfachen Übersicht aller Angebote wurde in den offenen Fragen mehrfach genannt. Manche Aussagen der Befragten zeigen außerdem eine fehlerhafte Einschätzung oder Verwechslung von Angeboten, was ebenfalls dazu beiträgt die Informationslage der Zielgruppe besser einschätzen zu können.

Für die Erhebung der Nutzung und der Bekanntheit (ET4) wurde für alle Angebote erhoben, ob die befragte Person das Angebot kennt und ob sie es in Anspruch nimmt. Außerdem wurde abgefragt, welche Angebote der Person bekannt sind, ohne diese vorher zu nennen. Damit kann ein bloßes Wiedererkennen vom Kennen der Angebote abgegrenzt werden. Die Untersuchung zeigt dabei eine besonders hohe Bekanntheit bei der KBC, dem Kita-Zuschuss, BuT und Wohngeld, während viele andere Angebote eher weniger bekannt.

Die Hürden (ET5) und Bedarfe (ET6) werden besonders in offenen Fragen ersichtlich. Dabei wurden die unterschiedlichen Nennungen geclustert. Als Hürden werden besonders der Informationsfluss und der bürokratische Aufwand betont. Dabei wurden von den Befragten auch konkrete Ideen geäußert wie der einer automatischen Verlängerung der KBC oder zumindest einer Erinnerung vor ihrem Ablauf. Ein Abgleich des Bedarfs mit der soziodemographischen Situation der Familien gibt dabei genauere Einblicke. So wird etwa der Mangel an Wohnraum von Familien mit mehr Kindern häufiger genannt und führt auch dazu, dass Kindern kein ausreichender Raum zum Lernen geboten werden kann.

Zusammenfassend ermöglicht die quantitative Methodik hier besonders detaillreiche Erkenntnisgewinne zur finanziellen Situation und Teilhabesituation, sowie zur Nutzung und Bekanntheit der Angebote. Zu den Hürden und Bedarfen konnten mittels offener Fragen Erkenntnisse gewonnen werden bis hin zu konkreten Vorschlägen. Rückschlüsse auf die Wirkung und Zielerreichung lassen sich hingegen auf Basis der Daten abgesehen von der messbaren Reichweite und qualitativen Kommentaren eher wenig bilden. Unter anderem dazu soll dagegen laut Studienautoren die qualitative Untersuchung und dabei besonders die Fokusgruppendiskussionenen Erkenntnisse liefern.

\subsection{Erkenntnisse der qualitativen Untersuchung}

\subsubsection{Literatur- und Datenanalyse}

Die Literatur- und Datenanalyse ermöglichte eine Einteilung in Handlungsfelder des Tübinger Präventionskonzepts gegen Kinderarmut. Dabei wurde erkannt, welche Angebote von welchen Trägern existieren, wieso sie ins Leben gerufen wurden und wie hoch ihre Inanspruchnahme ist. Von den Angeboten ausgehend konnten Rückschlüsse auf Bedürfnisse gezogen werden, die von den Angeboten befriedigt werden sollten. Außerdem wurden Maßnahmen erkannt, welche bereits getroffen wurden, um präventiv auf möglichen Schwächen der Angebote zu reagieren. Dazu seien die Thüringer Ansprechpartner (TAPs, TODO) genannt, welche einer mangelnden Bekanntheit der Angebote entgegenwirken. Neben den Handlungsfeldern wurde erkannt, dass das Präventionskonzept sowohl in unterschiedlichen Sozialräumen als auch in unterschiedlichen Lebensphasen unterstützende Angebote bietet. Bei den Erkenntnissen der Literatur- und Datenanalyse handelt es sich also weniger um Erkenntnisse zur Beantwortung des Forschungsinteresses als zur Entwicklung der Fragebögen zur quantitativen Untersuchung und der Leitfaden zur Fokusgruppendiskussion zur qualitativen Untersuchung.

\subsubsection{Fokusgruppen}

Fokusgruppendiskussionen wurden sowohl mit der Seite der haupt- und ehrenamtlichen Mitarbeitenden geführt, als auch mit der Zielgruppe der von Armut betroffenen oder bedrohten Familien in Tübingen. Hier findet also eine Triangulation unterschiedlicher Perspektiven statt. So ist es etwa für die Anspruchsberechtigten unverständlich, wieso Leistungen bei mehreren unterschiedlichen Stellen beantragt werden müssen, während Mitarbeitende die Hintergründe der Angebote und ihrer Finanzierungen besser nachvollziehen können. (S.109, TODO).

Zur Erhebung der Wirkung (ET1) konnten besonders die Aussagen der Zielgruppe zur Wahrnehmung der Angebote beitragen. So hebt diese die Angebote gegenüber anderen Städten hervor und auch die Mitarbeitenden nehmen ein positives Image einzelner Angebote wahr. Es wurden auch Wirkungen erkannt, die nicht unbedingt zur direkten Zielsetzung der Maßnahmen gehören, wie eine Steigerung der Sichtbarkeit von Armut (S.109, TODO) und der Sensibilität in den Institutionen und Vereinen (S.113, TODO). Erkenntnisse zur Zielerreichung (ET2) lassen sich aus dem Abgleich der Aussagen der Befragten mit den den Angeboten inhärenten Zielen erheben. So führt etwa der Peer-to-Peer-Ansatz des INET-Programms, bei dem Menschen als Ansprechpartner fungieren, die selbst eine Migrationsgeschichte haben, zur Wahrnehmung eines leichteren Kontaktaufbaus mit der Zielgruppe (S.131, TODO).

Aussagen über den Informationszugang (ET3), die Nutzung und Bekanntheit (ET4) sowie bestehende Hürden (ET5) werden von Befragten häufig zusammengefasst. Dabei werden insbesondere seitens der Mitarbeitenden unterschiedliche Theorien gebildet. Zum Beispiel, dass Stichwörter wie "KBC" erstmal bekannt werden müssen, bevor Personen im Internet danach suchen und die entsprechenden Ressourcen finden (S.139, TODO) oder dass Personen erst nach Informationen suchen, wenn durch einen persönlichen Kontakt darauf aufmerksam gemacht wurden (S.140, TODO). Die Äußerungen in den Fokusgruppen der von Armut betroffenen oder bedrohten Familien, bieten ebenfalls konkrete Rückschlüsse auf bestehende Hürden beim Informationsmaterial, wie die Hürde durch schwierige Sprache und unbekannte Fachbegriffe und Abkürzungen.

Der Erkenntnisgewinn über den Bedarf (ET6) ist durch die Mitarbeitenden und Familien sehr breit. So nennen Mitarbeitende etwa Optimierungspotential bei den Schulungen und Familien höhere Leistungen für bestimmte Bereiche wie Freiezeitangebote oder Schulbedarf (TODO). Dabei werden nicht nur Bedarfe gennant, sondern teilweise mit konkreten Verbesserungsvorschlägen verknüpft. So schlagen Mitarbeitende etwa vor, das die TAPs auch selbst lernen, Anträge auszufüllen, um bessere Hilfestellungen bieten zu können (TODO).

Die qualitative Untersuchung dieser Studie bietet somit einen breiten Erkenntnisgewinn über alle definierten Erkenntnistypen. Dieser geht teilweise über die Forschungsfragen hinaus und bietet eine Reihe an Theorien und Verbesserungsvorschlägen.

\subsection{Triangulation von Datenquellen, Methoden und theoretischen Perspektiven}

In den vorigen beiden Abschnitten wurden die Erkenntnisgewinne der qualitativen und der quantitativen Untersuchungen beschrieben. Auf dieser Basis soll nun analysiert werden, inwiefern die Kombination beider Paradigmen die Erkenntnisse bestätigt oder erweitert. Dabei wird die Triangulation auf drei Ebenen betrachtet: die Triangulation von Datenquellen, die Triangulation von Methoden und die Triangulation theoretischer Perspektiven. Dabei wird jeweils auf Umsetzung, gegenseitige Validierung und zusätzlichen Erkenntnisgewinn eingegangen.

\subsubsection{Triangulation von Datenquellen}

Die Tübinger Evaluationsstudie umfasst sowohl Primärdaten der quantitativen Umfrage und der qualitativen Forschungsgruppendiskussionen von Mitarbeitenden und Betroffenen als auch Sekundärdaten der Literatur- und Datenanalyse und amtlicher Statistiken zum SGB-II-Bezug. Die Verwendung amtlicher Statistiken zur Quotenwahl bei der Stichprobe der quanitativen Untersuchung erhöht die externe Validität. Die unterschiedlichen Erhebungen erreichen dabei unterschiedliche Personengruppen. So wurden etwa in der Online-Umfrage auch englisch-sprachige Personen und Personen, denen eine Anreise zu einer Fokusgruppe schwer möglich wäre, einbezogen. Neben den Fokusgruppen aus Eltern von Armut betroffener oder bedrohter Familien wurden auch Fokusgruppen von Mitarbeitenden gebildet. Damit lassen sich Aussagen der Eltern durch die Aussagen der Mitarbeitenden aus einer zweiten Perspektive validieren. So heben bereits die Studienautoren überschneidende und gegensetzliche Wahrnehmungen der beiden Gruppen hervor (TODO).

\subsubsection{Triangulation von Methoden}

Auf der Ebene der Methoden werden in der Studie die quantitative Methode der Befragung in Form einer Online-Umfrage und die qualitativen Methoden der Literatur- und Datenanalyse sowie der Fokusgruppendiskussionen eingesetzt. Im Zuge des sequenziellen, konvergenten Mixed-Methods-Design werden Erkenntnisse der einzelnen Methoden für die Konzeption nachfolgender Methoden verwendet (siehe Kapitel "Foschungsdesign der Studie", TODO).

Die statistischen Erkenntnisse der quantitativen Untersuchung zur Teilhabesituation der Kinder werden in den Fokusgruppendiskussionen der Eltern wiedergespiegelt. So zeigt die Deprivationsstudie, den höchsten Mangel in den Kategorien Wohnraum und Freizeitangebote (S. 71), die ebenfalls wiederholt von den Eltern in den Fokusgruppen benannt und vertieft ausgeführt werden (TODO). Ähnlich verhält es sich bei Reichweite und Bekanntheit der Angebote. Angebote, die in der Umfrage häufig genannt werden wie die KBC, werden auch von den Teilnehmenden der Fokusgruppen häufiger angesprochen.

\subsubsection{Triangulation theoretischer Perspektiven}

In der Tübinger Evaluationsstudie werden mindestens vier theoretische Perspektiven auf Kinderarmut angewendet. Erstens die relative Einkommensarmut, die, erweitert um die Frage nach dem Bezug von Sozialleistungen, als Kriterium für die Wahl der Stichprobe der quantitativen Untersuchung verwendet wurde. Zweitens die Deprivation, die auf den britischen Soziologen Peter Townsend zurückgeht (Handbuch Kinderarmut S.31). Verwendet wurde die Deprivationsanalyse in Anlehnung an den Fragebogen von World Vision (World Vision/Sabine Andresen et al. 2018) zur Untersuchung, inwieweit grundlegende Bedürfnisse für die Lebensqualität von Kindern erfüllt werden.

Drittens ist der Lebenslagenansatz implizit in den Ergebnissen der Fokusgruppendiskussion zu erkennen (siehe hierzu Weisser 1989/1956, TODO). Nach Wolfgang Voges (TODO) gibt es keine einheitlichen Indikatoren zur Messung der Lebenslage. Zentral ist der Handlungsspielraum der Personen, der sich in den Aussagen der Befragten zu ihrer Lebenssituation und ihren Herausforderungen zur gesellschaftlichen Teilhabe sowie in der Nutzung von Hilfsangeboten wiederfindet. Viertens wird in der quantitativen Untersuchung die soziale Segregation betrachtet (Katharina Knüttel und Volker Kersting). So wird etwa festgestellt, dass "sind diejenigen Befragten, deren finanzielle Lage schlecht oder sehr schlecht ist, anteilig häufiger im Sozialraum Nord (54 \%) oder in den Teilorten wohnhaft (53 \%)" (S. 62, TODO) sind.

Durch die Kombination der vier theoretischen Perspektiven entsteht ein breiteres Bild von den Verhältnisse in denen von Armut betroffene oder bedrohte Familien leben. Während die Konzepte der Einkommensarmut und Segregation Erkenntnisse über die Rahmenbedingungen liefern, vertiefen die Betrachung der Deprivation und der Lebenslage die Teilhabesituation und Lebensrealität der Familien.

\subsection{Beitrag zur Validität und Zuverlässigkeit der Ergebnisse}



- Quoten-Stichproben erhöhen externe Validität

Verbesserungen
- Perspektive der Kinder fehlt
- weitere theoretische Ansätze wie die sozialisationstheoretische Perspektive hilfreich
- explizitere Triangulation in Form eines Kapitels, der die quantitativen und qualitativen Erkenntnisse kombiniert
- Quoten der Stichproben tatsächlich erzielen
- 